%
% teil2.tex -- Beispiel-File für teil2 
%
% (c) 2020 Prof Dr Andreas Müller, Hochschule Rapperswil
%
% !TEX root = ../../buch.tex
% !TEX encoding = UTF-8
%
\section{Zusammenhang mit Möbius Transformation
\label{geradlinig:section:zusammenhang_mit_moebius_transformation}}
\kopfrechts{Zusammenhang mit Möbius Transformation}
Die Möbius Transformation ist die Abbildung von komplexen Funktionen
\begin{align*}
    z\mapsto\frac{az + b}{cz + d},
\end{align*}
wobei $a, b, c, d$ komplexe Zahlen sind und $ad - bc \neq 0$ gilt. 
Die Möbius Transformation wird in Kapitel~\ref{chapter:moebius}
noch detaillierter behandelt.
Hier beschränken wir uns lediglich auf den Elementartyp Inversion.

\subsection{Inversion}
Die Inversion wird beschrieben durch
\begin{align*}
    z\mapsto\frac{1}{z}
\end{align*}
